\section{The applications}
\label{app:generator}

\paragraph{Generator.} An application of $N$ task-chains is drawn chain by
chain from a seeded pseudo-random generator (xoshiro256** seeded by
splitmix64). The steps and their parameters are the following.
\begin{enumerate}
\item \emph{Structure} (WATERS 2015 rules \cite{kramer2015waters}): 1, 2 or
3 activation patterns with probabilities 70, 20 and 10\,\%; each pattern
has 2, 3, 4 or 5 tasks with probabilities 30, 40, 20 and 10\,\%. The
patterns of a task-chain have distinct periods, drawn with the shares of
\cref{tab:periods}.
\item \emph{Chain utilization}: log-uniform between 0.1 and 8\,\% of the
reference core; 2\,\% of the task-chains are heavy, uniform between 8 and
80\,\%.
\item \emph{Bindings}: 25\,\% of the task-chains are local to a random zone
(heavy ones never are). Of the others, 60\,\% start with a sensor task and
50\,\% end with an actuator task, each pinned to a random zone. On a
platform without zones, a random ECU that is not a central computer plays
the role of the zone. A
task-chain whose tasks would all be drivers keeps its last task free.
\item \emph{Task utilizations}: each driver takes 5\,\% of the chain
utilization; the computation tasks share the rest by UUniFast
\cite{bini2005measuring}, redrawn while a task exceeds 30\,\% of a core.
Then $C=\mathrm{round}(UT)$, at least \SI{1}{\micro s} and below $T$.
\item \emph{Messages}: 4--64 bytes (70\,\%), 65--1500 bytes (25\,\%), or
1501 bytes to 16\,kB (5\,\%).
\item \emph{Deadline}: 300\,\% of the sum of the periods of the patterns of
the task-chain, on its maximum reaction time.
\item \emph{Constraints}: memory per task log-uniform in 4--256\,KB.
Optionally, the experiments add: co-location of consecutive free tasks of
a pattern; zone-only task-chains; per-type WCET factors log-uniform in
$[(1+v)^{-1},1+v]$, lock-step equal to MCU; safety task-chains restricted to
lock-step cores; and replicas. A replica is a copy of a random non-heavy
task-chain, with its pinned tasks moved to zones that the original does
not use; the free tasks of each are kept out of the zones the other is
pinned to.
\end{enumerate}
The task-chains are finally sorted by decreasing utilization and numbered.
The task-chains of a software update (\cref{sec:lifecycle}) are drawn the
same way, from a seed derived from the seed of the application and the
number of the update, and numbered after the existing ones.

\begin{table}[ht]
\centering\small
\caption{Periods of the WATERS 2015 benchmark and their shares of the
runnables; the 15\,\% of angle-synchronous runnables are left out.}
\label{tab:periods}
\begin{tabular}{@{}lrrrrrrrrr@{}}
\toprule
period (ms) & 1 & 2 & 5 & 10 & 20 & 50 & 100 & 200 & 1000 \\
share (\%) & 3 & 2 & 2 & 25 & 25 & 3 & 20 & 1 & 4 \\
\bottomrule
\end{tabular}
\end{table}

\paragraph{The detection chain.} This illustrative ADAS chain is not taken
from a benchmark. It is one task-chain of nine tasks pinned to the central
computer, in the order GPS grabber (7/50), LIDAR grabber (12/50),
camera grabber (2/25), localization (28/50), detection (28/50), EKF
(6.5/10), fusion (25/50), planner (5/10) and control (4.5/10), where $C/T$
is in ms. The messages are 64, 16\,000, 16\,000, 256, 1024, 256, 1024, 256
and 64 bytes.

\section{The metaheuristic baselines}
\label{app:baselines}

Both baselines take the sub-chains of step 1 (\cref{sec:partitioning})
and map them onto a fixed set of \emph{candidate cores}: the initial cores
of the platform, then, for every ECU and core type, as many added cores as
the sub-chains allowed on that ECU whose lightest type it is there would
fill (their utilization on that type over $\varphi$, rounded up), within $M^{\max}_k$ and at most 64. A
mapping is admissible when it meets every constraint of
\cref{sec:problem}, schedulability included (the robust optimal priority
assignment of step 4 on every core). Both minimize the energy
\begin{equation}
E=M-\frac{0.9}{M}\sum_{p\ \text{used}}\Bigl(1-\frac{u_p}{\varphi}\Bigr)^2,
\label{eq:energy}
\end{equation}
where $M$ is the number of cores in use and $u_p$ their loads: fewer cores
first, then the free capacity gathered on few cores, the next ones to
empty. Since $E$ grows with the load, a sub-chain also prefers a
faster core. Neither runs the consolidation, and neither models
shared-memory contention. Their best mapping then gets the
deadline-driven local search of step 6, as in the heuristic, so that the
task-chains on time compare the searches. Both are deterministic for a
seed, and run on the default, slow-zone, constrained and domain
platforms of \cref{tab:exact}.

\paragraph{Simulated annealing} \cite{kirkpatrick1983optimization}. It
starts from a first fit of the sub-chains by decreasing utilization, each
to the first candidate core that takes it, the fastest cores first. A move
relocates one sub-chain to another candidate core (probability 0.7) or
exchanges two sub-chains of two cores (0.3). Half of the time, the moved
sub-chain is the one, of two drawn at random, on the less loaded core; the
target is a core in use with probability 0.8. A move that breaks a
constraint is not made, and one that raises the energy by $d$ is accepted
with probability $e^{-d/T}$, the temperature decreasing geometrically
from 0.03 to 0.0003 over the budget of moves. Every mapping visited is
admissible: a core that loses a sub-chain stays schedulable, since its
response times can only decrease.

\paragraph{Genetic algorithm.} An individual is an order of the
sub-chains, decoded by a packing heuristic, as in the hybrid genetic algorithms for
bin packing of Reeves \cite{reeves1996hybrid}. It is decoded by first fit: each sub-chain
goes to the first core in use that takes it (the fastest first, then in
the order they were opened), else opens the first candidate core that
does, the fastest first. The population has 40 individuals: the order by
decreasing utilization and random orders. Parents are chosen by binary
tournaments, recombined by order crossover (OX1, probability 0.9) and
mutated by a swap (probability 0.3); the two best individuals are kept.
